% Cover letter for a fresh submission, venue-neutral (Elsevier elsarticle manuscript).
% Replace the bracketed journal name before submitting. The previous version of this
% file was a revision letter addressed to the Journal of Statistical Planning and
% Inference (Ms. Ref. No. JSPI-D-26-00452), a round that closed with a rejection on
% presentation grounds; see report/REPORT.md for the submission history.
\documentclass[11pt,a4paper]{article}
\usepackage[margin=1in]{geometry}
\usepackage{parskip}
\usepackage{url}
\pagestyle{empty}

\begin{document}

\noindent
Quang-Vinh Dang\\
British University Vietnam\\
Hung Yen, Vietnam\\
\url{vinh.dq4@buv.edu.vn}

\vspace{1em}
\noindent
To the Editors,\\
\emph{[Journal name]}

\vspace{1em}

\noindent
Re: ``An Exact Decomposition of Paired Score Differences by Relabelling Within
Groups'' --- new submission

\vspace{1em}

Dear Editors,

Please consider the enclosed manuscript for publication in \emph{[Journal name]}.

\textbf{The problem, in one paragraph.} Two probabilistic classifiers are compared on a
test set and the newer one scores higher. If the labels depend partly on a grouping
variable that both models observe, the score difference sums two distinct improvements:
the new model may fit the label frequencies \emph{within} groups more closely, or it may
match its predictions to the individual cases they were made for more closely. These are
different achievements --- the first is skill a lookup table also supplies and that a
shift in label frequencies removes; the second is not --- and an aggregate score
difference cannot tell them apart. Neither can a significance test on that difference,
since both components are present in the same number.

\textbf{The contribution.} The paper decomposes the difference exactly, using no fitted
reference model. The construction is elementary: because both models emit a full
probability vector for every test case, each model's score can be recomputed against any
label, so one can score a case's predictions against a \emph{different} case's label from
the same group. That relabelling preserves the group and its label frequencies while
breaking the correspondence between a prediction and the case it was made for. The gain
that survives is the \emph{transported gain}; the gain that disappears is a
\emph{within-group covariance gain}. The resulting identity is exact in the observed
sample, for any score, with no held-out fold and no tuning parameter, at a cost of
$O(NK)$ in $N$ cases and $K$ labels.

The remainder is a U-statistic of order two, and it is not a new quantity: it equals the
within-group covariance between the two models' predicted-probability difference and the
label indicator, which is the object isolated by Yates's (1982) covariance decomposition
of a proper score, here evaluated for a \emph{pair} of forecasters instead of one. One
correction relative to earlier versions of this work bears directly on how the
contribution is framed, so I state it rather than leave it to be found: this remainder is
\emph{not} the resolution term of the classical reliability--resolution partition, as
earlier drafts claimed. Classical resolution depends on a forecast only through the
partition its level sets generate and is therefore invariant to any monotone rescaling of
forecast values; a covariance between those values is not. The two coincide only when both
models are calibrated within groups. The paper now states that relation precisely and
adopts confidence matching as part of the procedure rather than as a caveat. What the
classical decomposition does not supply in either reading is inference for the pairwise
contrast, and that is where the paper's technical content sits. Specifically, we prove: the covariance
identity for the whole Bregman-score family, so the quadratic and logarithmic scores are
named cases of one theorem; asymptotic normality of the dataset-level estimator with a
consistent cluster-robust sandwich variance; an exact reduction, at a trivial grouping
variable, to a clustered Diebold--Mariano/Giacomini--White test, which locates the
estimator precisely relative to comparative forecast evaluation; an exact $(n_c-1)/n_c$
attenuation of the in-sample plug-in, which is why no sample splitting is needed; an
explicit between-group covariance term for what happens when groups are merged; an exact
composition of the transported component into a mean-forecast-fit term minus a dispersion
term, which is what shows that component is not a measure of prior fit; and a
Cauchy--Schwarz sensitivity bound, in the tradition of Rosenbaum and Cinelli \& Hazlett,
for how far an unrecorded grouping variable could move the result. On that last point the
paper reports the bound against its own estimates rather than only in the abstract: for
one of the appendix studies an unrecorded covariate with an aggregate correlation of
$0.061$ would account for the whole reported gain, and I say so.

\textbf{Validation, and one application in detail.} Because the object is a measuring
instrument, the paper validates it before using it: signal recovery, null calibration,
coverage of the primary interval in the clustered small-group regime that stresses it,
and discriminant-validity designs establishing what the two components do \emph{not}
separate. One of those is a genuine limitation rather than a strength --- a monotone
recalibration moves score between the components while adding no information --- and the
paper responds with a calibration-matching protocol that is part of the procedure, plus a
demonstration that it works.

The main application audits two publicly released scene-graph checkpoints of a widely
cited debiasing method, evaluated in both modes with the original authors' code, so the
models are not ours and the claimed improvement is one the field already reports. Once
confidence is matched on held-out images and the audit runs at the object-class-pair
grouping, the proper-score difference is overwhelmingly transported under both the
quadratic and the logarithmic score, although the two scores disagree on whether any
covariance remainder survives ($-0.00006$, 95\% CI $[-0.00120,+0.00109]$ under the
quadratic score; $+0.04396$, CI $[+0.04079,+0.04714]$ under the log score) --- in either
reading at least an order of magnitude smaller than the transported shift. That
configuration qualifier appears in the paper wherever the claim does, because two of the
five configurations reported for this same pair do not support it. Since the method's
construction subtracts a context-only prediction, the finding characterises it rather than
indicting it, and the paper is explicit that the ten-point mean-recall movement and the
proper-score decomposition are in different units with no bridge supplied between them.

Three further studies --- our own relation predictors, a frozen-CLIP variant of them, and
long-tailed classification in images and in text --- are reported in an appendix, and
include a case in which a model that loses on every aggregate measure is shown to resolve
individual cases significantly better than the model it appears to lose to.

\textbf{Fit.} The central contribution is an estimand together with its exact
finite-sample and asymptotic inference theory. The benchmarks are where the instrument is
exercised and validated, not the claim: no state-of-the-art result is asserted, and none
is needed for the argument. I have written the paper so that the contribution is legible
before any notation appears --- a plain statement of the identity in the introduction, a
four-point worked example that exhibits the main properties proved later, and a table
fixing the paper's vocabulary against its standard statistical counterparts. The same
four points then exhibit a model-selection decision the decomposition gets right and the
aggregate score gets wrong: a candidate whose entire gain is transported becomes
\emph{worse} than the model it replaces under a shift in label frequencies, while one
whose gain is covariance is untouched.

All code, cached model outputs and scripts needed to reproduce every number, table and
figure are publicly available at \url{https://github.com/vinhqdang/WAGER}, including a
verification script that checks each reported value against the committed results files.

The manuscript is original, has not been published previously, and is not under
consideration for publication elsewhere. An earlier and substantially different version
was submitted to another journal and declined; that submission is closed. I am the sole
author and approve this submission. I declare no competing financial or personal
interests, and the work received no specific grant from any funding agency. In accordance
with journal policy, the manuscript includes a declaration of generative AI use in its
preparation.

Thank you for your consideration.

\vspace{1em}
\noindent
Sincerely,\\
Quang-Vinh Dang

\end{document}
